% ============================================================
% 3. METHOD
% ============================================================
% DRAFT v0.1 (2026-08-13). Solid draft, NOT final. Seams marked "FLAG:".
% DISCIPLINE: STRONG pathway is DESIGNED, NOT BUILT. Never say implemented/ran.
% Ashwin quote (22 Jul) woven in; accuracy guard applied (Parsa proposed the
% (b)-reading, Ashwin endorsed -- do NOT imply Ashwin proposed it end-to-end).
% ============================================================
\section{Method}
\contrib{lead: Parsa}

Our method is a frozen-representation probe of whether egocentric behavioral
signals carry action-anticipatory content, together with a controlled pathway
by which those signals are injected into the predictor. We first describe the
frozen-probe setup and evaluation protocol (\S\ref{subsec:method-setup}), then
the two conditioning pathways: a weak one that is built and verified, and a
stronger one that is deliberately specified but left unimplemented
(\S\ref{subsec:method-weak}--\ref{subsec:method-strong}). Last comes the
direct, signal-present readout used to measure the signal's information content
independent of any pathway (\S\ref{subsec:method-peva}).

% ------------------------------------------------------------
\subsection{Frozen-representation probing}
\label{subsec:method-setup}

We build on V-JEPA~2~\citep{assran2025vjepa2}, a self-supervised video model
whose encoder and latent predictor are pretrained to predict masked
spatiotemporal features rather than pixels. Throughout the anticipation probe, both the encoder and the
predictor are kept frozen; the only trained component is a lightweight attentive
probe that reads the predictor's output tokens and emits verb, noun, and action
logits for the EK100 anticipation task.

\paragraph{Backbone and projectors.}
Two predictor instantiations appear in this work, and we are explicit about which
carries which result. A ViT-L predictor (12 blocks, width 384) is used only for
the Phase-4 feasibility fine-tuning of \S\ref{subsec:protocol}; every
anticipation result reported below uses the ViT-G predictor (24 blocks, internal width
1024, projected back to the 1408-dimensional encoder space on exit), on a $224/16 = 14$ patch grid ($N = 196$ tokens per frame) over
$T = 8$ temporal positions. Behavioral signals enter through two single linear
projectors with no activation or normalisation: a gaze projector
($3\rightarrow1024$) and a hand projector ($12\rightarrow1024$), each paired with
a learned mask-token parameter used when the signal is invalid. Gaze carries one
validity flag; the two-hand signal carries two, combined by disjunction, so the
mask substitutes only when both hands are invalid and a one-hand-valid frame
passes its full twelve-dimensional vector (with the invalid half zero-filled).
Of the predictor's 24 blocks, the last six (blocks 18--23) are unfrozen during
conditioning. That six is the library default rather than an explicit argument,
confirmed empirically by the gradient guard (Appendix~\ref{sec:code}).

Freezing the backbone is a deliberate design choice rather than a shortcut: it
keeps the limited supervised capacity focused on the task's label space, limits
overfitting on rare classes, and makes the comparison between conditioning
pathways a comparison of \emph{inputs to a fixed predictor} rather than of two
separately fine-tuned models. The same frozen-probe recipe underlies the two
strongest published EK100 anticipation systems~\citep{chu2026jfaa,
wang2026tapjepa}, which lends the choice external support.

\paragraph{Masked evaluation is in-distribution.}
Several diagnostics in \S\ref{subsec:diagnostics} compare the predictor's
behaviour on the real signal against its behaviour when the signal is replaced
by a mask token. This masked condition is not out-of-distribution: the warm-start
predictor checkpoint was trained with a signal-dropout rate of $0.4$, so masked
behavioral slots were seen during training and the model has a defined response
to them. Comparisons against the mask token are therefore not measuring a distribution
shift. They do, however, conflate two effects --- the value of the signal's
content and the cost of withholding an input the model was trained to expect ---
and only a predictor trained without the signal at all can separate them. We report such a control in \S\ref{subsec:training-ablation}; it is null, which is what licenses
the content reading we take here. This dropout rate is the checkpoint's own training
setting: the warm-start predictor was fine-tuned within the project on
HD-EPIC~P01--P07 with P08 held out, for three epochs at a signal-dropout rate of
$0.4$, with the last six predictor blocks unfrozen.

% ------------------------------------------------------------
\subsection{The weak pathway: peer-token concatenation}
\label{subsec:method-weak}

The weak conditioning pathway is inherited from the project's earlier phases.
Gaze and hand are each encoded to a single token per frame and concatenated
alongside the visual tokens inside the predictor, mirroring the
action--state--visual token layout of the original action-conditioned predictor.
This widening is internal and transient: per frame the two behavioral tokens
take the predictor's working sequence from \num{1568} ($T\cdot N = 8\times196$)
to \num{1584} ($T\cdot(N{+}2) = 8\times198$), and they are stripped again before
the predictor returns, so its output is once more \num{1568} tokens. The
behavioral signal therefore conditions the prediction without appearing in the
predictor's output.

\paragraph{Bit-exact equivalence as a control.}
For the weak pathway to serve as a baseline, it must reproduce the established
Phase-5 evaluation route exactly; otherwise a later comparison would confound
the pathway with an incidental change in evaluation code. We verify an exact
match, with identical mean, standard deviation, and extrema, and a maximum
absolute difference of zero to eight decimal places over the full
$[1, 1568, 1408]$ output (Appendix~\ref{sec:code}). The weak pathway is thus a
faithful control, not an approximate one.

\paragraph{Guarding against silent nulls.}
Three mechanisms in this codebase can each produce the appearance that
``conditioning does not help'' without the conditioned modules ever being
trained, and all three fail silently: a \texttt{no\_grad} scope around the
predictor forward, a parameter-group misclassification that routes new modules
into a low-learning-rate group, and default weight initialisation in place of the
codebase's scheme. Because the last of these reproduces the very
weight-stagnation symptom the project set out to explain, we verify after a
single backward pass that every unfrozen module receives a non-zero gradient;
all ten watched modules pass and the frozen blocks correctly receive none
(Appendix~\ref{sec:code}). This guard is a precondition for any future
implementation of the stronger pathway.

% ------------------------------------------------------------
\subsection{The strong pathway: designed, not built}
\label{subsec:method-strong}

The predictor exposes a second conditioning mode behind the same switch,
reserved for a stronger pathway that \emph{adapts} mechanisms from
GazeQwen~\citep{pham2026gazeqwen} to our architecture: biasing the visual
tokens' keys with the behavioral signal, a bottleneck resampler that fuses gaze
and hand, and a learned per-block gate over the fine-tuned predictor blocks. We
stress that this is an adaptation rather than a reproduction: GazeQwen injects
its signal externally, as a resampler residual added onto language-model decoder
layers, and performs no in-place modification of a vision transformer's
attention; the in-attention key-biasing here is our own design choice, informed
by their mechanisms but structurally distinct from their implementation.
\textbf{This pathway is deliberately left unimplemented.} Construction of the
switch succeeds, but its forward pass raises rather than silently running a
partial pathway (Appendix~\ref{sec:code}); no result in this paper is produced by
the strong pathway, and it is never trained or evaluated.

The ordering is intentional. Following supervision, we adopted the
mechanism-adaptation reading of GazeQwen, adapting its key-biasing and resampler
mechanisms into our own predictor, rather than reproducing its pipeline wholesale
(its Qwen2.5-VL host, its benchmark, and live gaze), which was judged to offer
little.\footnote{Our supervisor's guidance was explicit
on this point: to ``go with option (b),'' the mechanism adaptation, because
``there's nothing interesting that would come out of'' a one-to-one replication
of the GazeQwen paper, and to ``perform the direct test as planned'' before
committing to implementation.} A stronger pathway is worth building only if a
stronger pathway could help; \S\ref{subsec:diagnostics} tests exactly that before
any implementation cost is incurred, and \S\ref{subsec:b4} records the resulting
decision to retire it at the design stage. The design stays on record, and because its zero-initialised gate makes it
runnable as a falsifier, the choice can be revisited if the sample-size
constraint of \S\ref{subsec:ddn} is relieved.

% ------------------------------------------------------------
\subsection{Direct readout of signal content}
\label{subsec:method-peva}

Independent of any injection pathway, we measure how much anticipatory
information the behavioral signal carries by reading it out directly, adapting
the whole-body-conditioned prediction protocol of
PEVA~\citep{bai2025peva} to a signal-present readout on HD-EPIC~P01. We
distinguish two questions the readout can answer. The first, an average-error
comparison (PEVA-1), asks whether the predictor's output is on average closer to
the target when the real signal is present than when it is masked, which tests
detectability. The second, a discrimination comparison (PEVA-2), asks whether the
signal opens a usable margin between correct and incorrect candidates, which tests
whether any detectable effect is large enough to change a decision. As
\S\ref{subsec:diagnostics} reports, the signal is detectable under the first test
and negligible under the second. That decoupling, a real but decision-irrelevant
effect, is what the diagnostics turn on.